\documentclass[12pt,a4paper]{report}
\usepackage[utf8]{inputenc}
\usepackage{amsmath}
\usepackage{amsfonts}
\usepackage{amssymb}
\usepackage{amsthm}
\usepackage{hyperref}

\usepackage{multicol}
\usepackage{fancyhdr}
\usepackage[inline]{enumitem}
\usepackage{tikz}
\usepackage{tikz-cd}
\usetikzlibrary{calc}
\usetikzlibrary{shapes.geometric}
\usepackage[margin=0.5in]{geometry}
\usepackage{xcolor}
\usepackage{ esint }

\hypersetup{
    colorlinks=true,
    linkcolor=blue,
    filecolor=magenta,      
    urlcolor=cyan,
    pdftitle={Real Analysis},
    pdfpagemode=FullScreen,
    }

%\urlstyle{same}

\newcommand{\CLASSNAME}{Comprehensive Exams -- Real Analysis}
\newcommand{\STUDENTNAME}{Paul Carmody}
\newcommand{\ASSIGNMENT}{NOTES }
\newcommand{\DUEDATE}{April 2026}
\newcommand{\SEMESTER}{Spring 2026}
\newcommand{\SCHEDULE}{TBD}
\newcommand{\ROOM}{Crawford 332}

\newcommand{\MMN}{M_{m\times n}}
\newcommand{\FF}{\mathcal{F}}

\pagestyle{fancy}
\fancyhf{}
\chead{ \fancyplain{}{\CLASSNAME} }
%\chead{ \fancyplain{}{\STUDENTNAME} }
\rhead{\thepage}
\fancyfoot[C]{-- \thepage --}
\input{../MyMacros}

\newcommand{\RED}[1]{\textcolor{red}{#1}}
\newcommand{\BLUE}[1]{\textcolor{blue}{#1}}
\newcommand{\FFF}{\mathcal{F}}

\begin{document}

\begin{center}
	\Large{\CLASSNAME -- \SEMESTER} 
\end{center}
\begin{center}
	\STUDENTNAME \\
	\ASSIGNMENT -- \DUEDATE\\
\end{center} 

\section*{Real Analysis -- Overview}

Here’s a structured overview of \DEFINE{Real Analysis}, organized the way it’s typically developed mathematically. I’ll emphasize the core ideas, key definitions, and fundamental theorems rather than just listing topics.

---

\begin{description}
\item \DEFINE{I. Foundations: The Real Number System}

	\begin{description}
    \item \textbf{1. Construction and Properties of $\R$}

		\begin{itemize}
        \item Ordered field axioms
        \item Completeness axiom (crucial!):
		\end{itemize}

  * Every nonempty set bounded above has a least upper bound (supremum)

    \item \textbf{2. Consequences of Completeness}

		\begin{itemize}
        \item Archimedean property
        \item Density of $\Q$ in $\R$
        \item Existence of limits
		\end{itemize}

	\end{description}
---

\item \DEFINE{II. Sequences and Limits}

	\begin{description}
    \item \textbf{1. Convergence of Sequences}

		\begin{itemize}
        \item Definition:
  \begin{align*}
  x_n \to x \iff \forall \varepsilon > 0, \exists N \text{ such that } n \ge N \Rightarrow |x_n - x| < \varepsilon
  \end{align*}
		\end{itemize}

    \item \textbf{2. Key Theorems}

		\begin{itemize}
        \item Uniqueness of limits
        \item Algebra of limits
        \item Squeeze Theorem
        \item Monotone Convergence Theorem
        \item Bolzano–Weierstrass Theorem (subsequence convergence)
        \item Cauchy criterion:
		\end{itemize}

  * Sequence converges $\iff$ it is Cauchy (uses completeness)

	\end{description}
---

\item \DEFINE{III. Series}

	\begin{description}
    \item \textbf{1. Infinite Series}

		\begin{itemize}
        \item Definition of convergence via partial sums
		\end{itemize}

    \item \textbf{2. Tests for Convergence}

		\begin{itemize}
        \item Comparison test
        \item Ratio test
        \item Root test
        \item Integral test
        \item Alternating series test
		\end{itemize}

    \item \textbf{3. Absolute vs Conditional Convergence}

		\begin{itemize}
        \item Absolute convergence $\to$ convergence
        \item Rearrangement theorem (Riemann)
		\end{itemize}

	\end{description}
---

\item \DEFINE{IV. Topology of $\R^n$}

	\begin{description}
    \item \textbf{1. Open and Closed Sets}

		\begin{itemize}
        \item Open sets ($\varepsilon$-balls)
        \item Closed sets (contain limit points)
		\end{itemize}

    \item \textbf{2. Compactness}

		\begin{itemize}
        \item Definition: every open cover has a finite subcover
        \item In $\R^n$:
		\end{itemize}

  * Heine–Borel Theorem: compact $\iff$ closed and bounded

    \item \textbf{3. Connectedness}

		\begin{itemize}
        \item Intervals as connected sets
        \item Intermediate Value property
		\end{itemize}

	\end{description}
---

\item \DEFINE{V. Limits and Continuity of Functions}

	\begin{description}
    \item \textbf{1. Limits of Functions}

		\begin{itemize}
        \item $\varepsilon-\delta$ definition
		\end{itemize}

    \item \textbf{2. Continuity}

		\begin{itemize}
        \item Equivalent formulations:
		\end{itemize}

  * Sequential continuity
  * $\varepsilon-\delta$ definition

    \item \textbf{3. Key Theorems}

		\begin{itemize}
        \item Intermediate Value Theorem (IVT)
        \item Extreme Value Theorem (EVT)
        \item Uniform continuity (vs pointwise continuity)
		\end{itemize}

	\end{description}
---

\item \DEFINE{VI. Differentiation}

	\begin{description}
    \item \textbf{1. Derivative}

		\begin{itemize}
        \item Definition:
  \begin{align*}
  f'(x) = \lim_{h \to 0} \frac{f(x+h) - f(x)}{h}
  \end{align*}
		\end{itemize}

    \item \textbf{2. Main Results}

		\begin{itemize}
        \item Mean Value Theorem (MVT)
        \item Rolle’s Theorem
        \item Taylor’s Theorem
		\end{itemize}

    \item \textbf{3. Consequences}

		\begin{itemize}
        \item Monotonicity criteria
        \item Convexity
        \item L’Hôpital’s Rule
		\end{itemize}

	\end{description}
---

\item \DEFINE{VII. Riemann Integration}

	\begin{description}
    \item \textbf{1. Definition}

		\begin{itemize}
        \item Upper and lower sums
        \item Integrability $\iff$ upper = lower integral
		\end{itemize}

    \item \textbf{2. Key Results}

		\begin{itemize}
        \item Continuous $\to$ integrable
        \item Monotone $\to$ integrable
        \item Fundamental Theorem of Calculus:
		\end{itemize}

  * Part I: differentiation of integrals
  * Part II: evaluation via antiderivatives

	\end{description}
---

\item \DEFINE{VIII. Sequences and Series of Functions}

	\begin{description}
    \item \textbf{1. Types of Convergence}

		\begin{itemize}
        \item Pointwise convergence
        \item Uniform convergence
		\end{itemize}

    \item \textbf{2. Key Theorems}

		\begin{itemize}
        \item Uniform convergence preserves:

  * Continuity
  * Integrability
  * Differentiability (with conditions)
		\end{itemize}

    \item \textbf{3. Important Tests}

		\begin{itemize}
        \item Weierstrass M-test
        \item Arzelà–Ascoli Theorem
		\end{itemize}

	\end{description}
---

\item \DEFINE{IX. Metric Spaces (Generalization)}

	\begin{description}
    \item \textbf{1. Metric Spaces}

		\begin{itemize}
        \item Distance function $d(x,y)$
		\end{itemize}

    \item \textbf{2. General Concepts}

		\begin{itemize}
        \item Convergence
        \item Continuity
        \item Compactness
        \item Completeness
		\end{itemize}

    \item \textbf{3. Important Results}

		\begin{itemize}
        \item Contraction Mapping Theorem (Banach Fixed Point)
		\end{itemize}

	\end{description}
---

\item \DEFINE{X. Measure Theory and Lebesgue Integration (Advanced Real Analysis)}

	\begin{description}
    \item \textbf{1. Motivation}

		\begin{itemize}
        \item Riemann integral limitations
		\end{itemize}

    \item \textbf{2. Measure}

		\begin{itemize}
        \item $\sigma$-algebras
        \item Lebesgue measure
		\end{itemize}

    \item \textbf{3. Lebesgue Integral}

		\begin{itemize}
        \item Integrate more general functions
		\end{itemize}

    \item \textbf{4. Convergence Theorems}

		\begin{itemize}
        \item Monotone Convergence Theorem (MCT)
        \item Dominated Convergence Theorem (DCT)
        \item Fatou’s Lemma
		\end{itemize}

	\end{description}
---

\item \DEFINE{XI. Functional Analysis Connections}

		\begin{itemize}
        \item Normed spaces, Banach spaces
        \item $L^p$ spaces
        \item Weak convergence
		\end{itemize}

---

\item \DEFINE{Core Themes to Keep in Mind}

Real Analysis revolves around a few deep ideas:

		\begin{itemize}
        \item \DEFINE{Completeness} → guarantees limits exist
        \item \DEFINE{Approximation} → sequences, series, integrals
        \item \DEFINE{Structure of sets} → topology
        \item \DEFINE{Limiting processes} → continuity, differentiation, integration
        \item \DEFINE{Control of convergence} → uniform vs pointwise
		\end{itemize}

---
\end{description}

If you want, I can tailor this into:

		\begin{itemize}
        \item a \DEFINE{course roadmap} (week-by-week),
        \item a \DEFINE{theorem checklist for exams}, or
        \item a \DEFINE{dependency map} showing which results rely on which foundations.
		\end{itemize}

\newpage
\section*{Dependency Map}

Here’s a \DEFINE{dependency map of Real Analysis}, showing how the major concepts and theorems build on one another. Think of it as a layered structure: each level relies on the ones below it.

---

\begin{description}
\item \DEFINE{0. Logical Foundations}

		\begin{itemize}
        \item Set theory (sets, functions, relations)
        \item Logic (quantifiers, implications, proof techniques)
		\end{itemize}

$\downarrow$ everything depends on this

---

\item \DEFINE{1. Real Numbers ($\R$) — The Core Foundation}

		\begin{itemize}
        \item Ordered field structure
        \item \DEFINE{Completeness axiom (supremum property)}
		\end{itemize}

This is the *most important root node*.

$\downarrow$ gives rise to:

		\begin{itemize}
        \item Archimedean property
        \item Density of rationals
        \item Nested interval property
		\end{itemize}

---

\item \DEFINE{2. Sequences}

Built directly from completeness.

	\begin{itemize}
    \item Depends on:

		\begin{itemize}
        \item $\R$ (especially completeness)
		\end{itemize}

    \item Leads to:

		\begin{itemize}
        \item Definition of limits
        \item Cauchy sequences
        \item Subsequence arguments
		\end{itemize}

    \item Key results:

		\begin{itemize}
        \item Cauchy $\iff$ convergent (\DEFINE{uses completeness})
        \item Bolzano–Weierstrass Theorem
        \item Monotone Convergence Theorem
		\end{itemize}
	\end{itemize}

$\downarrow$ feeds into:

---

\item \DEFINE{3. Series}

	\begin{itemize}
    \item Depends on:

		\begin{itemize}
        \item Sequences (via partial sums)
		\end{itemize}

    \item Leads to:

		\begin{itemize}
        \item Convergence tests
        \item Absolute vs conditional convergence
		\end{itemize}
	\end{itemize}

$\downarrow$ parallel development with:

---

\item \DEFINE{4. Topology of $\R^n$}

	\begin{itemize}
    \item Depends on:

		\begin{itemize}
        \item Distance (metric from $\R$)
		\end{itemize}

    \item Core concepts:

		\begin{itemize}
        \item Open/closed sets
        \item Limit points
        \item Compactness
        \item Connectedness
		\end{itemize}

    \item Key theorem:

		\begin{itemize}
        \item \DEFINE{Heine–Borel} (uses completeness indirectly)
		\end{itemize}

	\end{itemize}
$\downarrow$ feeds into:

---

\item \DEFINE{5. Limits of Functions}

	\begin{itemize}
    \item Depends on:

		\begin{itemize}
        \item Sequences OR topology

Two equivalent frameworks:

        \item $\varepsilon-\delta$ definition
        \item Sequential definition
		\end{itemize}

	\end{itemize}
$\downarrow$ leads to:

---

\item \DEFINE{6. Continuity}

	\begin{itemize}
    \item Depends on:

		\begin{itemize}
        \item Limits of functions
        \item Topology
		\end{itemize}

    \item Key results:

		\begin{itemize}
        \item Intermediate Value Theorem (uses connectedness)
        \item Extreme Value Theorem (uses compactness)
		\end{itemize}

	\end{itemize}
$\downarrow$ feeds into:

---

\item \DEFINE{7. Differentiation}

	\begin{itemize}
    \item Depends on:

		\begin{itemize}
        \item Limits
		\end{itemize}

    \item Core results:

		\begin{itemize}
        \item Mean Value Theorem (\DEFINE{depends on EVT + IVT $\to$ topology!})
        \item Taylor’s Theorem
		\end{itemize}

	\end{itemize}
$\downarrow$ feeds into:

---

\item \DEFINE{8. Riemann Integration}

	\begin{itemize}
    \item Depends on:

		\begin{itemize}
        \item Limits
        \item Supremum/infimum ($\to$ completeness)
		\end{itemize}

    \item Key results:

		\begin{itemize}
        \item Integrability criteria
        \item Fundamental Theorem of Calculus
		\end{itemize}

\DEFINE{ Important dependency:}

		\begin{itemize}
        \item FTC uses differentiation results (MVT structure)
		\end{itemize}

	\end{itemize}
$\downarrow$ feeds into:

---

\item \DEFINE{9. Sequences of Functions}

	\begin{itemize}
    \item Depends on:

		\begin{itemize}
        \item Sequences
        \item Functions
        \item Continuity/integration
		\end{itemize}

    \item Key ideas:

		\begin{itemize}
        \item Pointwise vs uniform convergence
		\end{itemize}

    \item Key results:

		\begin{itemize}
        \item Uniform convergence preserves continuity/integration
		\end{itemize}

	\end{itemize}
$\downarrow$ leads to:

---

\item \DEFINE{10. Metric Spaces (Generalization Layer)}

Abstracts everything above.

	\begin{itemize}
    \item Depends on:

		\begin{itemize}
        \item Understanding of $\R$ structure
		\end{itemize}

    \item Generalizes:

		\begin{itemize}
        \item Sequences $\to$ convergence
        \item Open/closed sets $\to$ topology
        \item Continuity
		\end{itemize}

    \item Key theorem:

		\begin{itemize}
        \item Contraction Mapping Theorem (uses completeness)
		\end{itemize}
	\end{itemize}

$\downarrow$ feeds into:

---

\item \DEFINE{11. Measure Theory \& Lebesgue Integration}

	\begin{itemize}
    \item Depends on:

		\begin{itemize}
        \item Set theory
        \item Limits
        \item Sequences of functions
		\end{itemize}

    \item Key shift:

		\begin{itemize}
        \item Replace “interval-based” integration with “measure-based”
		\end{itemize}

    \item Core theorems:

		\begin{itemize}
        \item Dominated Convergence Theorem
        \item Monotone Convergence Theorem
		\end{itemize}
	\end{itemize}

---

\item \DEFINE{Big Picture Dependency Flow}

Here’s the structure in compact form:

```
%Set Theory / Logic
%        ↓
%   Real Numbers (Completeness)
%        ↓
%     Sequences
%        ↓
%   ┌───────────────┐
%   ↓               ↓
% Series        Topology ($\R$ⁿ)
%                    ↓
%              Function Limits
%                    ↓
%                Continuity
%                    ↓
%              Differentiation
%                    ↓
%             Riemann Integration
%                    ↓
%        Sequences of Functions
%                    ↓
%        Metric Spaces (Abstract)
%                    ↓
%     Measure Theory (Lebesgue)
%```

---

\item \DEFINE{Key Insight (What Really Drives Everything)}

If you compress the whole map to its essence:

		\begin{itemize}
        \item \DEFINE{Completeness of $\R$} $\to$ makes limits work
        \item \DEFINE{Limits} $\to$ define continuity, derivatives, integrals
        \item \DEFINE{Topology} $\to$ explains *why* key theorems (IVT, EVT) hold
        \item \DEFINE{Uniform control} $\to$ determines when limits behave nicely
		\end{itemize}

---
\end{description}

If you want, I can turn this into a \DEFINE{visual diagram (graph-style)} or give you a \DEFINE{minimal set of “core theorems”} that generate almost everything else.

\newpage
\section*{11. Measure Theory -- Overview}


Here’s a structured overview of \DEFINE{Measure Theory}, focusing on the key ideas, objects, and theorems that make the subject work.

---
\begin{description}

\item \DEFINE{1. Motivation}

Measure theory generalizes the idea of:
		\begin{itemize}
        \item \DEFINE{length} (intervals),
        \item \DEFINE{area} (regions),
        \item \DEFINE{volume} (higher dimensions),
		\end{itemize}
so that we can:
		\begin{itemize}
        \item integrate more functions than Riemann allows,
        \item handle limits of functions cleanly,
        \item work rigorously with probability.
		\end{itemize}

---

\item \DEFINE{2. $\sigma$-Algebras (Measurable Structure)}

	\begin{description}
    \item \textbf{Definition}

A \DEFINE{$\sigma$-algebra} $\mathcal{F} \subseteq \mathcal{P}(X)$ is a collection of subsets such that:

		\begin{description}
        \item 1. $X \in \mathcal{F}$
        \item 2. Closed under complements: $A \in \mathcal{F} \Rightarrow A^c \in \mathcal{F}$
        \item 3. Closed under countable unions:
   \begin{align*} 
       A_1, A_2, \dots \in \mathcal{F} \Rightarrow \bigcup_{n=1}^\infty A_n \in \mathcal{F}
   \end{align*} 
   
		\end{description}

    \item \textbf{Interpretation}

		\begin{description}
        \item These are the sets to which we are allowed to assign “size.”
        \item Prevents pathological constructions.
		\end{description}

    \item \textbf{Key Example}

* \DEFINE{Borel $\sigma$-algebra} on $\R$: generated by open sets
	\end{description}

---

\item \DEFINE{3. Measures}

	\begin{description}
    \item \textbf{Definition}

A \DEFINE{measure} is a function:
\begin{align*}
\mu: \mathcal{F} \to [0,\infty]
\end{align*}
such that:
		\begin{description}
        \item 1. $\mu(\emptyset) = 0$
        \item 2. \DEFINE{Countable additivity}:
     \begin{align*}
         \mu\left(\bigcup_{n=1}^\infty A_n\right) = \sum_{n=1}^\infty \mu(A_n) \quad \text{(disjoint sets)}
     \end{align*}
		\end{description}

    \item \textbf{Examples}

		\begin{itemize}
        \item Length on $\R$ $\to$ \DEFINE{Lebesgue measure}
        \item Counting measure
        \item Probability measures $(\mu(X)=1)$
		\end{itemize}
	\end{description}

---

\item \DEFINE{4. Measurable Functions}

	\begin{description}
    \item \textbf{Definition}

A function $f: X \to \mathbb{R}$ is \DEFINE{measurable} if:
\begin{align*}
\{x : f(x) > a\} \in \mathcal{F} \quad \forall a \in \mathbb{R}
\end{align*}

    \item \textbf{Intuition}

		\begin{itemize}
        \item Measurable functions are those compatible with the measure structure.
        \item Continuous functions are measurable, but measurable functions are much more general.
\BLUE{
		\item The range of $f$ is measurable, (i.e., a member of a $\sigma$-algebra).
}
		\end{itemize}
	\end{description}

---

\item \DEFINE{5. Lebesgue Integration}

	\begin{description}
    \item \textbf{Step 1: Simple Functions}

* Finite linear combinations of indicator functions:
    \begin{align*}
        s(x) = \sum a_i \chi_{A_i}(x)
    \end{align*}

    \item \textbf{Step 2: Nonnegative Functions}

* Approximate from below by simple functions:
    \begin{align*}
        \int f\, d\mu = \sup \left\{ \int s\, d\mu : s \le f \right\}
    \end{align*}\BLUE{Decompose $f$ via its domain into a sequence of simple functions $s_i$ then integrate over the domain of each.  Keep in mind that this could easily be an infinite sum.  In Riemann we summed over the range using $\Delta x_i f(x_i)$ here we are using almost the same thing but $s_i(x)$ over each domain. Even more important is $X$ can be any multi-dimensional space whereas $\Delta x_i$ becomes limited.}

    \item \textbf{Step 3: General Functions}

* Decompose:
    \begin{align*}
        f = f^+ - f^-
    \end{align*}\BLUE{This process only really works if $f>0$, thus, for those functions that don't fit into this set, we must use $f^-$.}

---

    \item \textbf{Key Idea}

		\begin{itemize}
        \item \DEFINE{Riemann}: slice the domain
        \item \DEFINE{Lebesgue}: slice the range
		\end{itemize}

\BLUE{
		\item \textbf{Observations: } Up until this level, functions are largely regarded as an object that accepts a parameter and returns another object.  This emphasizes the idea that one input generates one output.  However, at this point, we focus on the domain as an intricate part of what a function is.  Thus, a simple function really IS about as simple a definition as we can from this perspective.  A simple function is a defined domain, $X$, and the function, $s: X \to \R$, is some constant.  That is $s(x)=c$ whenever $x \in X$ and zero otherwise.  The question now arises, what are the properties of $X$ (continuous, measurable, connected, complete, dense, $\dots$)?  These then determine the behavior of $s$.
}

	\end{description}

This makes Lebesgue integration much more flexible.

---

\item \DEFINE{6. Convergence Theorems (Core Power)}

These are the heart of measure theory.

	\begin{description}
    \item \textbf{1. Monotone Convergence Theorem (MCT)}

If $f_n \uparrow f$, then:
\begin{align*}
\int f_n \to \int f
\end{align*}

---

    \item \textbf{2. Fatou’s Lemma}

\begin{align*}
\int \liminf f_n \le \liminf \int f_n
\end{align*}\BLUE{The integral of the limit of the infimum is the `floor' of the sequence of functions.  Whereas the limit of the infimum of integrals will always be higher.}

---

    \item \textbf{3. Dominated Convergence Theorem (DCT)}

If:
		\begin{itemize}
        \item $f_n \to f$ pointwise
        \item $|f_n| \le g$, with $g \in L^1$
		\end{itemize}

Then:
\begin{align*}
\int f_n \to \int f
\end{align*}\BLUE{If $g$ exists it must be finite.  This is essentially an upper bound for $f_n$ and since $f_n \to f$ then $f$ must also be finite and hence $f\in L^1$ (set of integrable functions).}

---

\item Why these matter

They allow you to:
		\begin{itemize}
        \item interchange limits and integrals,
        \item control sequences of functions,
        \item justify calculations rigorously.
		\end{itemize}
	\end{description}

---

\item \DEFINE{7. $L^p$ Spaces}

	\begin{description}
    \item \textbf{Definition}

\begin{align*}
L^p(X) = \left\{ f : \int |f|^p < \infty \right\}
\end{align*}

    \item \textbf{Special cases}

		\begin{itemize}
        \item $L^1$: integrable functions
        \item $L^2$: Hilbert space (inner product structure)
		\end{itemize}

    \item \textbf{Key inequalities}

		\begin{itemize}
        \item Hölder’s inequality \BLUE{Think Cauchy-Schwartz
		\begin{align*}
			\sum |a_ib_i|  &\le \|a\|_{L^p} \|b\|_{L^q}
		\end{align*}
		}
        \item Minkowski inequality \BLUE{Think Triangular Inequality
		\begin{align*}
			\PAREN{\sum_n |a_i-b_i|^p}^{1/p} &\le \PAREN{\sum_n |a_i|^p}^{1/p} + 			\PAREN{\sum_n |b_i|^p}^{1/p} \\
			\PAREN{\int_n |a_i-b_i|^p}^{1/p} &\le \PAREN{\int_n |a_i|^p}^{1/p} + 			\PAREN{\int_n |b_i|^p}^{1/p}
		\end{align*}
        }

		\end{itemize}
	\end{description}

---

\item \DEFINE{8. Null Sets and “Almost Everywhere”}

	\begin{description}
    \item \textbf{Null sets}

* Sets of measure zero

    \item \textbf{Almost everywhere (a.e.)}

* Property holds except on a null set

    \item \textbf{Importance}

* Functions that differ only on null sets are considered the same in $L^p$
	\end{description}

---

\item \DEFINE{9. Product Measures and Fubini’s Theorem}

	\begin{description}
    \item \textbf{Product measure}

* Measure on $X \times Y$

    \item \textbf{Fubini’s Theorem}

\begin{align*}
    \int f(x,y)\,d(x,y) = \int \left(\int f(x,y)\,dx\right) dy
\end{align*}

    \item \textbf{Importance}

* Justifies iterated integrals
	\end{description}

---

\item \DEFINE{10. Radon–Nikodym Theorem}

If $\nu \ll \mu$, then:
\begin{align*}
    \nu(A) = \int_A f \, d\mu
\end{align*}

* $f = \frac{d\nu}{d\mu}$ is the \DEFINE{Radon–Nikodym derivative}

\textbf{Interpretation}

* Generalized ``density"

---

\item \DEFINE{11. Big Conceptual Themes}

	\begin{description}
    \item \textbf{1. Generalization of size}

* From intervals $\to$ arbitrary sets

    \item \textbf{2. Limit control}

* Convergence theorems dominate the subject

    \item \textbf{3. Flexibility}

* Handles discontinuous, highly irregular functions

    \item \textbf{4. Equivalence up to null sets}

* Exact equality replaced by “almost everywhere”
	\end{description}

---

\item \DEFINE{12. Dependency Flow (Compressed)}

%```id="7k2mql"
%$\sigma$-algebras
%     ↓
%Measures
%     ↓
%Measurable functions
%     ↓
%Lebesgue integral
%     ↓
%Convergence theorems
%     ↓
%Lp spaces / Functional analysis
%```

---
\end{description}

\DEFINE{If you want to go deeper}

I can walk you through:

		\begin{itemize}
        \item a \DEFINE{full construction of Lebesgue measure} (Carathéodory approach),
        \item \DEFINE{proofs of MCT/DCT} (these are central),
        \item or how measure theory connects directly to \DEFINE{probability theory} and \DEFINE{PDEs}.
		\end{itemize}

Just tell me where you want to zoom in.

\end{document}
